\subsection{The symmetric restriction map and its homological limitation}

\begin{proposition}[P-F71: finite normalization under symmetric restriction]
\label{prop:symmetricrestriction}
Let $n\ge4$, $V=\mathbb C^n$, $W=V^{\otimes3}$,
$S=\operatorname{Sym}^3V\subset W$, and $g=\gcd(n,3)$.
Give the following rings the displayed stretch grading:
\[
 R=\bigoplus_{s\ge0}\mathbb C[W]_{ns}^{SL(V)^3},\qquad
 B=\bigoplus_{s\ge0}\mathbb C[S]_{ns}^{SL(V)}.
                                                               \tag{P-F71a}
\]
Thus $R=R_n^{\mathrm{ten}}$ of P-F60, whereas $B$ is the
$g$th Veronese, regraded, of the actual cubic ring of P-F19.
In particular
\[
 \dim R_s=k((s^n),(s^n),(s^n)),\qquad
 \dim B_s=f_n(gs).                                             \tag{P-F71b}
\]
Let $\rho:R\longrightarrow B$ be polynomial restriction,
and let $A=\rho(R)$.  This is an actual graded algebra
homomorphism.  The inclusion $A\subset B$ is finite and
birational; $B$ is its normalization.  The induced finite
map $\operatorname{Spec}B\to\operatorname{Spec}A$ is
bijective on complex points.

More precisely, set
\[
 G'=\{D\in GL(V):(\det D)^3=1\}.
                                                               \tag{P-F71c}
\]
For any two symmetric tensors, equivalence under $SL(V)^3$
holds if and only if they are in the same diagonal $G'$-orbit.
The $SL(V)^3$ nullcone intersects $S$ in exactly the
diagonal $SL(V)$ nullcone, and $B=\mathbb C[S]^{G'}$ with
the grading in (P-F71a).

There is a nonnegative integer $c_n$, not bounded here,
and a nonzero homogeneous element $h\in A_{c_n}$ such that
$hB\subset A$.  Consequently, for every $s\ge0$,
\[
 f_n(gs)\ \le\ \dim A_{s+c_n}\ \le\
              k(((s+c_n)^n),((s+c_n)^n),((s+c_n)^n)).
                                                               \tag{P-F71d}
\]
No assertion that $A=B$, or that $c_n$ is polynomially
bounded, is included.  In fact the inclusion is strict
on the infinite family identified below.
\end{proposition}

\begin{proof}
\emph{Determinants in the equivalence--congruence theorem.}
Belitskii--Sergeichuk, \emph{Congruence of multilinear
forms}, Theorem~2(a), printed p.~3, and Corollary~8,
printed p.~9, prove that complex symmetric multilinear
forms are equivalent precisely when they are congruent:
\url{https://arxiv.org/pdf/0710.0834}.
We need an additional determinant bookkeeping fact from
their proof, printed pp.~7--8, equation~(12), which we
verify explicitly.

Write an equivalence between symmetric trilinear forms as
$F(u_1,u_2,u_3)=G(\phi_1u_1,\phi_2u_2,\phi_3u_3)$.
Symmetry gives the same identity for every permutation
of the three maps.  In the source induction, suppose
$\phi_1=\cdots=\phi_t=\phi$, where $t=1$ or $2$.
Its construction supplies a $G$-selfadjoint invertible
map $\eta$ with
\[
 \eta^{t+1}=(\phi\phi_{t+1}^{-1})^{-1},
 \qquad
 \psi_1=\cdots=\psi_{t+1}=\eta\phi,
 \quad \psi_j=\phi_j\quad(j>t+1).
\]
The construction uses polynomial roots on the generalized
eigenspaces; selfadjointness means that a map may be moved
from one argument to another.  It therefore preserves
the equivalence.  It also preserves the determinant product,
since
\[
 (\det\eta)^{t+1}(\det\phi)^{t+1}
   =\det(\phi\phi_{t+1}^{-1})^{-1}(\det\phi)^{t+1}
   =(\det\phi)^t\det\phi_{t+1}.
                                                               \tag{P-F71e}
\]
After the two steps the congruence map $D$ thus satisfies
$(\det D)^3=\prod_{j=1}^3\det\phi_j$.  Apply this to
the multilinear forms on $V^*$ associated with the tensors;
duals or inverses do not alter the condition that this
product is one.  An $SL(V)^3$ equivalence therefore gives
a diagonal $G'$ congruence.

Conversely, if $\delta=\det D$ and $\delta^3=1$, choose
$a,b\in\mathbb C^*$ with $a^n=b^n=\delta^{-1}$ and put
$c=(ab)^{-1}$.  Then $c^n=\delta^2=\delta^{-1}$.
All three maps $aD,bD,cD$ have determinant one and their
tensor product has the same action as $D^{\otimes3}$.
This proves the exact orbit assertion, with no assumption
of smoothness, a scalar centroid, or a generic tensor.

\emph{The grading and nullcones.}
The centers of the three special linear factors show
that ordinary tensor-invariant degrees are divisible by
$n$.  A homogeneous diagonal $SL(V)$ invariant of
ordinary degree $d$ has $n\mid3d$ and transforms under
$GL(V)$ by the character $\det^{3d/n}$ (or its inverse
with the coordinate-action convention).  It is invariant
under $G'$ precisely when $3\mid3d/n$, that is, when
$n\mid d$.  Thus $\mathbb C[S]^{G'}$ is exactly $B$.
Since the full cubic ring uses ordinary degree $n/g$ as
one stretch unit, this also proves the Veronese assertion
and (P-F71b), including the multiplication and dual
conventions already verified in P-F19 and P-F60.

Use a unitary structure on $V$ and the induced structure
on $W$.  The three reduced Hermitian matrices of a
symmetric tensor are equal.  Its diagonal $SU(V)$ moment
map is the sum of the three marginal moment maps, so
it vanishes if and only if each marginal moment map
vanishes.  By Kempf--Ness, every closed diagonal
$SL(V)$ orbit therefore remains closed for $SL(V)^3$.
The precise closed-orbit criterion used here is Losev,
\emph{The Kempf--Ness theorem and invariant theory},
Theorem~4, printed p.~1 and proof on p.~2:
\url{https://arxiv.org/pdf/math/0605756}.

Certainly a diagonally unstable symmetric tensor is
unstable for the larger group.  Conversely, suppose
$F\in S$ is diagonally semistable.  Its diagonal orbit
closure contains a nonzero closed diagonal orbit,
represented by $F_0$.  This orbit is also closed for
$SL(V)^3$ by the preceding paragraph.  If $F$ were
$SL(V)^3$-unstable, every positive-degree invariant of
that group would vanish at $F$ and hence at $F_0$.
This would put the nonzero closed orbit of $F_0$ in
the zero fiber, a contradiction.  Here we use the
affine reductive quotient property that each fiber
has a unique closed orbit, as in Luna, \emph{Slices
\'etales}, Section~I.1, printed p.~84:
\url{https://www.numdam.org/item/10.24033/msmf.110.pdf}.
This proves equality of nullcones.  The finite extension
$SL(V)\subset G'$ does not change the nullcone or the
closed-orbit property.

\emph{Finiteness and normalization.}
Both invariant rings are finitely generated.  Normality
of $B$ also follows directly: an element of
$\operatorname{Frac}(B)$ integral over $B$ is integral
over the polynomial ring $\mathbb C[S]$, hence belongs
to that ring; being a fraction of invariants it is
$G'$-invariant and therefore belongs to $B$.

Let $A_+$ and $B_+$ denote the positive-degree ideals.
The equality of nullcones and surjectivity of the affine
quotient $S\to\operatorname{Spec}B$ show that
\[
                       \sqrt{A_+B}=B_+.                       \tag{P-F71f}
\]
Indeed a point annihilating $A_+$ has a symmetric
representative in the larger-group nullcone, and hence
has the vertex as its image under the diagonal quotient.
It follows that $B/A_+B$ is finite dimensional.  Lift a
homogeneous vector-space basis of this quotient to $B$.
These finitely many lifts generate $B$ as an $A$-module:
in each degree, subtract their scalar combinations and
write the remainder in $A_+B$; induction on degree
expresses its lower-degree $B$ factors using the same
lifts.  Thus $A\subset B$ is finite.

To prove injectivity on closed points, represent two
points of $\operatorname{Spec}B$ by closed $G'$ orbits.
Their diagonal special linear orbits are closed, because
$G'/SL(V)$ is finite.  They are therefore closed also
under $SL(V)^3$.  Equal values of $A$ mean equal values
of all larger-group invariants, so these two closed
larger-group orbits coincide.  The determinant-preserving
orbit assertion above then identifies the two $G'$
orbits.  Finiteness gives surjectivity on closed points.

In characteristic zero a finite dominant map of integral
varieties which is injective on closed points is birational:
otherwise the separable function-field extension has
degree greater than one and, after removing the
discriminant locus of a primitive element, a closed
fiber contains more than one geometric point.
Consequently $A$ and $B$ have the same fraction field.
Since $B$ is finite and normal, it is exactly the integral
closure of $A$ in that field.  This establishes the
claimed normalization, without replacing normalization
by an unproved isomorphism.

Finally choose finitely many $A$-module generators of
$B$.  Since all belong to $\operatorname{Frac}(A)$,
a common nonzero denominator gives an element
$h\in A$ with $hB\subset A$.  The conductor ideal
$\{a\in A:aB\subset A\}$ is graded: this can be checked
on homogeneous module generators by taking homogeneous
components of each product.  Hence $h$ may be chosen
homogeneous, of some degree $c_n\ge0$.  Multiplication
by $h$ is injective because $B$ is a domain.  Applying
it to $B_s$ and using the surjection
$R_{s+c_n}\to A_{s+c_n}$ proves (P-F71d).
\end{proof}

\begin{corollary}[The finite-projective-dimension shortcut fails]
For every $n\ge4$, the finite graded $R$-module $B$
in P-F71 has infinite graded projective dimension.
More strongly,
\[
                         \frac{H_B(t)}{H_R(t)}
                    \notin\mathbb C[t,t^{-1}].                \tag{P-F71g}
\]
Thus there is no finite graded free resolution giving
a Laurent-polynomial Hilbert-series transfer from the
tensor invariant ring to this symmetric restriction module.
\end{corollary}

\begin{proof}
For a nonzero rational function, write $\deg_\infty$
for numerator degree minus denominator degree after
cancellation.  We first record the exact primary degree
bound being used.  Knop--Littelmann,
\emph{Der Grad erzeugender Funktionen von Invariantenringen},
Math. Z. 196 (1987), Section~1, equations~(1)--(2),
printed p.~211, write the ordinary invariant Hilbert
series with $q=-\deg_\infty H$.
Section~2, printed p.~212, the paragraph beginning
``In [9], Beweis von Korollar 4'', identifies
$\dim W-q$ as the degree of a polynomial covariant;
in particular $q\le\dim W$:
\url{https://doi.org/10.1007/BF01163656}.
The primary author-uploaded scan is
\url{https://www.researchgate.net/profile/Friedrich-Knop/publication/225947227_Der_Grad_erzeugender_Funktionen_von_Invariantenringen_The_degree_of_generating_functions_of_rings_of_invariants/links/0c9605262b324671cb000000/Der-Grad-erzeugender-Funktionen-von-Invariantenringen-The-degree-of-generating-functions-of-rings-of-invariants.pdf}.
The acting effective quotient of $SL(V)^3$ is connected
semisimple over $\mathbb C$, so this bound applies.
Ordinary degrees in $R$ are $n$ times the stretch degrees;
therefore
\[
                        \deg_\infty H_R\ge-n^2.                \tag{P-F71h}
\]

The actual function $\dim B_s=f_n(gs)$ is a nonzero
global stretching quasipolynomial.  GCT VI,
Theorem~1.6.1(b), printed p.~21, explicitly gives strict
properness of its Hilbert series:
\url{https://arxiv.org/pdf/0704.0229}.
Equivalently, apply the elementary global-quasipolynomial
properness argument of P-F9 to this Veronese sequence.
Thus, in the inherited integer grading,
\[
                        \deg_\infty H_B\le-1.                 \tag{P-F71i}
\]
There is no inference here from a statement about only
eventual quasipolynomials.  The canonical identification
in P-F19 and the stretching by $g$ are precisely what
allow this global statement.

Put $m=\binom{n+2}{3}$.  The closed orbits with finite
effective stabilizer in P-F24 and P-F60 give
\[
 \dim R=n^3-3n^2+3,\qquad
 \dim B=m-n^2+1,
 \qquad
 \Delta_n=\dim R-\dim B=n^3-2n^2+2-m.                          \tag{P-F71j}
\]
Suppose the quotient in (P-F71g) were a Laurent
polynomial $P(t)$.  Both Hilbert series are regular
at zero with value one, so $P$ is an ordinary polynomial
with $P(0)=1$.  Equations (P-F71h)--(P-F71i) give
$\deg P\le n^2-1$.  On the other hand, the orders of
the Hilbert-series poles at $t=1$ equal the respective
Krull dimensions.  Their positive leading coefficients
give a zero of $P$ of exact order $\Delta_n$.
This is impossible, since for $n\ge5$,
\[
 \Delta_n-(n^2-1)
  =\frac{(n-5)(5n^2+4n+18)+108}{6}>0.                          \tag{P-F71k}
\]
For $n=4$, use the already verified Salmon series
(P-F35b).  Here $g=1$ and
\[
 \deg_\infty H_B=25-(2+4+6+8+10)=-5,
 \qquad \Delta_4=14.
                                                               \tag{P-F71l}
\]
Together with $\deg_\infty H_R\ge-16$, a polynomial
quotient would have degree at most $11$ and a zero of
order $14$ at one, again impossible.  This proves
(P-F71g) for every $n\ge4$.

A finite graded free resolution would give
$H_B/H_R=\sum_{i,j}(-1)^i\beta_{ij}t^{d_{ij}}$, a Laurent
polynomial, by taking its Hilbert-series Euler characteristic.
Finite graded projective resolutions give the same
conclusion: finitely generated graded projective modules
over the connected positive graded algebra $R$ are graded
free, by lifting a basis modulo $R_+$ and applying graded
Nakayama to the resulting split surjection.  Thus such a
resolution cannot exist.
\end{proof}

\begin{corollary}[Explicit nonsurjectivity and a nonsplit normalization]
For each integer $r\ge2$, set
\[
                    n=\binom{3r-1}{2}.
                                                               \tag{P-F71m}
\]
Then the restriction map in P-F71 is not surjective:
\[
                    \dim B_r=1,\qquad \dim R_r=0.
                                                               \tag{P-F71n}
\]
Thus $A\subsetneq B$, the ring $A$ is not normal,
and its normalization inclusion has no $A$-linear
retraction.  Every nonzero homogeneous conductor
element $h\in A_c$ satisfies
\[
                            c\ge\lceil\sqrt n\rceil.
                                                               \tag{P-F71o}
\]
The first instance is $n=10$, $r=2$.
\end{corollary}

\begin{proof}
We have $n=(3r-1)(3r-2)/2\equiv1\pmod3$, so $g=1$.
Apply the exact sign-twist exterior identity (P-F32a)
with $k=3r$ and $d=nr$:
\[
 B_r\ \text{has dimension}\quad
 \operatorname{mult}_{\det^n}
       \bigwedge^{nr}\!\left(\bigwedge^3\mathbb C^{3r}\right).
                                                               \tag{P-F71p}
\]
Here $\dim\bigwedge^3\mathbb C^{3r}=nr$, and the determinant
of this exterior-cube representation is
$\det^{\binom{3r-1}{2}}=\det^n$: each diagonal variable
occurs in exactly $\binom{3r-1}{2}$ basis weights.
Thus (P-F71p) is the multiplicity of a one-dimensional
character in itself and equals one.  This is the top
exterior-power case of P-F32; its sign-twist identity
and primary locators are fully proved there.

On the tensor side, transposing the first two Specht
partitions tensors both by the sign character, whose
square is trivial.  Therefore
\[
 k((r^n),(r^n),(r^n))
          =k((n^r),(n^r),(r^n))=0.
                                                               \tag{P-F71q}
\]
The vanishing is the Schur--Weyl length bound (P-F58e):
the first two partitions have length $r$, whereas the
third has length $n>r^2$, since
\[
                  n-r^2=\frac{(7r-2)(r-1)}2>0.
\]
Its primary tensor-product identity is Manivel,
\emph{On rectangular Kronecker coefficients}, Section~2,
equation~(1), printed pp.~1--2:
\url{https://arxiv.org/pdf/0907.3351}.
This proves (P-F71n) and hence $A\subsetneq B$.
Since $B$ is the normalization of $A$ in their common
fraction field, $A$ is not normal.

If an $A$-linear retraction $p:B\to A$ existed,
take any $b\in B$ and a nonzero denominator $a\in A$
with $ab\in A$.  Then
$a\,p(b)=p(ab)=ab$.  Cancellation in their common
fraction field gives $p(b)=b$, forcing every element
of $B$ to belong to $A$, a contradiction.

Finally the same length argument shows $R_c=0$
whenever $0<c<\sqrt n$.  A nonzero conductor element
of degree zero would be a scalar, and would imply
$B\subset A$, already excluded.  Since $A_c$ is a
quotient of $R_c$, every possible positive conductor
degree is at least $\lceil\sqrt n\rceil$, proving
(P-F71o).  This lower bound is polynomial; no larger
or effective upper bound is asserted.
\end{proof}

\paragraph{What this exact restriction dictionary does and does not transfer.}
The map in P-F71 respects multiplication and has a finite
normalization with a shifted coefficient inequality.
These conclusions close the orbit, grading and finiteness
questions for this map.  The infinite family (P-F71m)
explicitly refutes universal surjectivity and an
$A$-linear splitting of the normalization.  No upper
bound on its conductor or equality between the two
stretching functions has been established.
Even the additional finite-resolution criterion that
would force $H_B=H_R P$ fails for all $n\ge4$.
Consequently this restriction result supplies no
ordinary Hilbert-pole transfer and no positivity-index
or saturation-index transfer to the rectangular
Kronecker ray.  The multiple fixed components from
P-F64--65 remain a separate issue.
